% =====================================================================
% 摘要与关键词
% ---------------------------------------------------------------------
% 用法：放入主文档 \begin{document} 之后、正文之前：
%     % =====================================================================
% 摘要与关键词
% ---------------------------------------------------------------------
% 用法：放入主文档 \begin{document} 之后、正文之前：
%     % =====================================================================
% 摘要与关键词
% ---------------------------------------------------------------------
% 用法：放入主文档 \begin{document} 之后、正文之前：
%     % =====================================================================
% 摘要与关键词
% ---------------------------------------------------------------------
% 用法：放入主文档 \begin{document} 之后、正文之前：
%     \input{abstract}
% 适用 ctexart / ctexbook。若需英文摘要，把下面的 english 部分放在正文末尾。
% 依赖：无额外宏包（ctex 自带 abstract 环境）
% =====================================================================

\begin{abstract}
\noindent
央行《货币政策执行报告》是市场识别政策意图的主要书面渠道。已有研究确认报告文本
情绪对股市整体收益具有解释力，并指出真正起作用的是面向未来的政策指引而非宏观回顾；
但这些结论多以宽基指数为对象，较少讨论不同制度阶段是否一致。

本文以 2001 年第一季度至 2026 年第二季度共 102 期报告为样本，采用中文金融情感词典
与情感单元法，分别度量政策指引段与宏观回顾段的情绪；以 16 个申万一级行业的等权组合
构造事件窗收益与波动指标，考察报告发布后一日与五日内的市场反应。

估计结果显示，政策指引段情绪对发布后的收益方向没有稳健影响，却显著降低事件日的市场
波动：情绪每提高一个标准差，日内振幅下降 15 至 28 个基点（取决于是否包含金融危机年份），
控制发布前波动状态后为 15 至 19 个基点；该结果在 28 行业口径、500 次置换检验与
政策操作控制下均保持稳健，而宏观回顾段情绪无此效应。分时段来看，波动效应集中于
2014 年以前；2014 年以后情绪与发布后收益转为负相关，但这一方向性结果对危机年份高度
敏感，本文仅将其作为提示性证据。

上述发现意味着，央行书面沟通的作用首先体现在收敛市场分歧，而非改变价格方向，且这种
作用随货币政策框架的调整而变化。若把不同时期、不同文本段落视为同质样本处理，方向相反
的两类效应会相互抵消，容易得出沟通无效的错误判断。
\end{abstract}

\vspace{0.5em}
\noindent\textbf{关键词：}货币政策执行报告；文本情绪；市场波动；结构断点；事件研究

\vspace{0.3em}
\noindent\textbf{JEL 分类号：}E52；E58；G14

% ---------------------------------------------------------------------
% 英文摘要（按需保留；若期刊或赛事不要求可删去）
% ---------------------------------------------------------------------
\clearpage
\renewcommand{\abstractname}{Abstract}
\begin{abstract}
\noindent
The Monetary Policy Report (MPR) published by the People's Bank of China is the
primary written channel through which policy intentions are communicated to the
market. Existing studies have established that the tone of the MPR explains
aggregate stock market returns, and that the forward-looking policy guidance
rather than the macroeconomic review drives this reaction. Most of this evidence,
however, rests on broad market indices, and whether it holds across institutional
periods has received little attention. Using 102 quarterly reports from 2001Q1 to
2026Q2 and a Chinese financial sentiment dictionary, we measure the tone of the two
segments separately. Market reactions are measured with event-window returns and
volatility constructed from an equal-weighted portfolio of 16 Shenwan level-1
industries.

We find no robust effect of policy guidance tone on the direction of post-release
returns, but a significant reduction in announcement-day volatility: a one
standard deviation increase in tone lowers the intraday range by 15 to 28 basis
points, depending on whether crisis years are included, and by 15 to 19 basis
points after controlling for pre-announcement volatility. The result survives an
alternative 28-industry market definition, a 500-draw permutation test, and
controls for reserve requirement ratio cuts, LPR cuts and reverse repo rate cuts,
while the tone of the macroeconomic review shows no such effect. The volatility
effect is concentrated before 2014; after 2014 the relation between tone and
returns turns negative, but this directional result is highly sensitive to crisis
years and is therefore reported as suggestive evidence.

These findings suggest that written central bank communication primarily reduces
disagreement rather than moving prices, and that its effect is not constant but
varies with the monetary policy framework. Treating heterogeneous periods and text
segments as a homogeneous sample may cancel out effects of opposite sign and lead
to the misleading conclusion that communication does not matter.
\end{abstract}

\vspace{0.5em}
\noindent\textbf{Keywords:} Monetary Policy Report; textual sentiment; market volatility;
structural break; event study

% ---------------------------------------------------------------------
% 写作依据（仅供团队核对，不进入排版）
%   "面向未来的政策指引而非宏观回顾"        → 姜富伟等（2021）的核心结论
%   前瞻性/回顾性分段的处理思路              → 王辉等（2023）；蒋海等（2025）
%   文本度量对词典与设定的敏感性（稳健性）    → 林建浩、孙乐轩（2025）
%   "信息效应"而非单纯冲击的分析框架          → 王熙、郑梦圆（2021）
%   行业事件窗与行业面板做法                  → 游庭玮（2026）；王春芝（2022）；周亮（2018）
% ---------------------------------------------------------------------

% 适用 ctexart / ctexbook。若需英文摘要，把下面的 english 部分放在正文末尾。
% 依赖：无额外宏包（ctex 自带 abstract 环境）
% =====================================================================

\begin{abstract}
\noindent
央行《货币政策执行报告》是市场识别政策意图的主要书面渠道。已有研究确认报告文本
情绪对股市整体收益具有解释力，并指出真正起作用的是面向未来的政策指引而非宏观回顾；
但这些结论多以宽基指数为对象，较少讨论不同制度阶段是否一致。

本文以 2001 年第一季度至 2026 年第二季度共 102 期报告为样本，采用中文金融情感词典
与情感单元法，分别度量政策指引段与宏观回顾段的情绪；以 16 个申万一级行业的等权组合
构造事件窗收益与波动指标，考察报告发布后一日与五日内的市场反应。

估计结果显示，政策指引段情绪对发布后的收益方向没有稳健影响，却显著降低事件日的市场
波动：情绪每提高一个标准差，日内振幅下降 15 至 28 个基点（取决于是否包含金融危机年份），
控制发布前波动状态后为 15 至 19 个基点；该结果在 28 行业口径、500 次置换检验与
政策操作控制下均保持稳健，而宏观回顾段情绪无此效应。分时段来看，波动效应集中于
2014 年以前；2014 年以后情绪与发布后收益转为负相关，但这一方向性结果对危机年份高度
敏感，本文仅将其作为提示性证据。

上述发现意味着，央行书面沟通的作用首先体现在收敛市场分歧，而非改变价格方向，且这种
作用随货币政策框架的调整而变化。若把不同时期、不同文本段落视为同质样本处理，方向相反
的两类效应会相互抵消，容易得出沟通无效的错误判断。
\end{abstract}

\vspace{0.5em}
\noindent\textbf{关键词：}货币政策执行报告；文本情绪；市场波动；结构断点；事件研究

\vspace{0.3em}
\noindent\textbf{JEL 分类号：}E52；E58；G14

% ---------------------------------------------------------------------
% 英文摘要（按需保留；若期刊或赛事不要求可删去）
% ---------------------------------------------------------------------
\clearpage
\renewcommand{\abstractname}{Abstract}
\begin{abstract}
\noindent
The Monetary Policy Report (MPR) published by the People's Bank of China is the
primary written channel through which policy intentions are communicated to the
market. Existing studies have established that the tone of the MPR explains
aggregate stock market returns, and that the forward-looking policy guidance
rather than the macroeconomic review drives this reaction. Most of this evidence,
however, rests on broad market indices, and whether it holds across institutional
periods has received little attention. Using 102 quarterly reports from 2001Q1 to
2026Q2 and a Chinese financial sentiment dictionary, we measure the tone of the two
segments separately. Market reactions are measured with event-window returns and
volatility constructed from an equal-weighted portfolio of 16 Shenwan level-1
industries.

We find no robust effect of policy guidance tone on the direction of post-release
returns, but a significant reduction in announcement-day volatility: a one
standard deviation increase in tone lowers the intraday range by 15 to 28 basis
points, depending on whether crisis years are included, and by 15 to 19 basis
points after controlling for pre-announcement volatility. The result survives an
alternative 28-industry market definition, a 500-draw permutation test, and
controls for reserve requirement ratio cuts, LPR cuts and reverse repo rate cuts,
while the tone of the macroeconomic review shows no such effect. The volatility
effect is concentrated before 2014; after 2014 the relation between tone and
returns turns negative, but this directional result is highly sensitive to crisis
years and is therefore reported as suggestive evidence.

These findings suggest that written central bank communication primarily reduces
disagreement rather than moving prices, and that its effect is not constant but
varies with the monetary policy framework. Treating heterogeneous periods and text
segments as a homogeneous sample may cancel out effects of opposite sign and lead
to the misleading conclusion that communication does not matter.
\end{abstract}

\vspace{0.5em}
\noindent\textbf{Keywords:} Monetary Policy Report; textual sentiment; market volatility;
structural break; event study

% ---------------------------------------------------------------------
% 写作依据（仅供团队核对，不进入排版）
%   "面向未来的政策指引而非宏观回顾"        → 姜富伟等（2021）的核心结论
%   前瞻性/回顾性分段的处理思路              → 王辉等（2023）；蒋海等（2025）
%   文本度量对词典与设定的敏感性（稳健性）    → 林建浩、孙乐轩（2025）
%   "信息效应"而非单纯冲击的分析框架          → 王熙、郑梦圆（2021）
%   行业事件窗与行业面板做法                  → 游庭玮（2026）；王春芝（2022）；周亮（2018）
% ---------------------------------------------------------------------

% 适用 ctexart / ctexbook。若需英文摘要，把下面的 english 部分放在正文末尾。
% 依赖：无额外宏包（ctex 自带 abstract 环境）
% =====================================================================

\begin{abstract}
\noindent
央行《货币政策执行报告》是市场识别政策意图的主要书面渠道。已有研究确认报告文本
情绪对股市整体收益具有解释力，并指出真正起作用的是面向未来的政策指引而非宏观回顾；
但这些结论多以宽基指数为对象，较少讨论不同制度阶段是否一致。

本文以 2001 年第一季度至 2026 年第二季度共 102 期报告为样本，采用中文金融情感词典
与情感单元法，分别度量政策指引段与宏观回顾段的情绪；以 16 个申万一级行业的等权组合
构造事件窗收益与波动指标，考察报告发布后一日与五日内的市场反应。

估计结果显示，政策指引段情绪对发布后的收益方向没有稳健影响，却显著降低事件日的市场
波动：情绪每提高一个标准差，日内振幅下降 15 至 28 个基点（取决于是否包含金融危机年份），
控制发布前波动状态后为 15 至 19 个基点；该结果在 28 行业口径、500 次置换检验与
政策操作控制下均保持稳健，而宏观回顾段情绪无此效应。分时段来看，波动效应集中于
2014 年以前；2014 年以后情绪与发布后收益转为负相关，但这一方向性结果对危机年份高度
敏感，本文仅将其作为提示性证据。

上述发现意味着，央行书面沟通的作用首先体现在收敛市场分歧，而非改变价格方向，且这种
作用随货币政策框架的调整而变化。若把不同时期、不同文本段落视为同质样本处理，方向相反
的两类效应会相互抵消，容易得出沟通无效的错误判断。
\end{abstract}

\vspace{0.5em}
\noindent\textbf{关键词：}货币政策执行报告；文本情绪；市场波动；结构断点；事件研究

\vspace{0.3em}
\noindent\textbf{JEL 分类号：}E52；E58；G14

% ---------------------------------------------------------------------
% 英文摘要（按需保留；若期刊或赛事不要求可删去）
% ---------------------------------------------------------------------
\clearpage
\renewcommand{\abstractname}{Abstract}
\begin{abstract}
\noindent
The Monetary Policy Report (MPR) published by the People's Bank of China is the
primary written channel through which policy intentions are communicated to the
market. Existing studies have established that the tone of the MPR explains
aggregate stock market returns, and that the forward-looking policy guidance
rather than the macroeconomic review drives this reaction. Most of this evidence,
however, rests on broad market indices, and whether it holds across institutional
periods has received little attention. Using 102 quarterly reports from 2001Q1 to
2026Q2 and a Chinese financial sentiment dictionary, we measure the tone of the two
segments separately. Market reactions are measured with event-window returns and
volatility constructed from an equal-weighted portfolio of 16 Shenwan level-1
industries.

We find no robust effect of policy guidance tone on the direction of post-release
returns, but a significant reduction in announcement-day volatility: a one
standard deviation increase in tone lowers the intraday range by 15 to 28 basis
points, depending on whether crisis years are included, and by 15 to 19 basis
points after controlling for pre-announcement volatility. The result survives an
alternative 28-industry market definition, a 500-draw permutation test, and
controls for reserve requirement ratio cuts, LPR cuts and reverse repo rate cuts,
while the tone of the macroeconomic review shows no such effect. The volatility
effect is concentrated before 2014; after 2014 the relation between tone and
returns turns negative, but this directional result is highly sensitive to crisis
years and is therefore reported as suggestive evidence.

These findings suggest that written central bank communication primarily reduces
disagreement rather than moving prices, and that its effect is not constant but
varies with the monetary policy framework. Treating heterogeneous periods and text
segments as a homogeneous sample may cancel out effects of opposite sign and lead
to the misleading conclusion that communication does not matter.
\end{abstract}

\vspace{0.5em}
\noindent\textbf{Keywords:} Monetary Policy Report; textual sentiment; market volatility;
structural break; event study

% ---------------------------------------------------------------------
% 写作依据（仅供团队核对，不进入排版）
%   "面向未来的政策指引而非宏观回顾"        → 姜富伟等（2021）的核心结论
%   前瞻性/回顾性分段的处理思路              → 王辉等（2023）；蒋海等（2025）
%   文本度量对词典与设定的敏感性（稳健性）    → 林建浩、孙乐轩（2025）
%   "信息效应"而非单纯冲击的分析框架          → 王熙、郑梦圆（2021）
%   行业事件窗与行业面板做法                  → 游庭玮（2026）；王春芝（2022）；周亮（2018）
% ---------------------------------------------------------------------

% 适用 ctexart / ctexbook。若需英文摘要，把下面的 english 部分放在正文末尾。
% 依赖：无额外宏包（ctex 自带 abstract 环境）
% =====================================================================

\begin{abstract}
\noindent
央行《货币政策执行报告》是市场识别政策意图的主要书面渠道。已有研究确认报告文本
情绪对股市整体收益具有解释力，并指出真正起作用的是面向未来的政策指引而非宏观回顾；
但这些结论多以宽基指数为对象，较少讨论不同制度阶段是否一致。

本文以 2001 年第一季度至 2026 年第二季度共 102 期报告为样本，采用中文金融情感词典
与情感单元法，分别度量政策指引段与宏观回顾段的情绪；以 16 个申万一级行业的等权组合
构造事件窗收益与波动指标，考察报告发布后一日与五日内的市场反应。

估计结果显示，政策指引段情绪对发布后的收益方向没有稳健影响，却显著降低事件日的市场
波动：情绪每提高一个标准差，日内振幅下降 15 至 28 个基点（取决于是否包含金融危机年份），
控制发布前波动状态后为 15 至 19 个基点；该结果在 28 行业口径、500 次置换检验与
政策操作控制下均保持稳健，而宏观回顾段情绪无此效应。分时段来看，波动效应集中于
2014 年以前；2014 年以后情绪与发布后收益转为负相关，但这一方向性结果对危机年份高度
敏感，本文仅将其作为提示性证据。

上述发现意味着，央行书面沟通的作用首先体现在收敛市场分歧，而非改变价格方向，且这种
作用随货币政策框架的调整而变化。若把不同时期、不同文本段落视为同质样本处理，方向相反
的两类效应会相互抵消，容易得出沟通无效的错误判断。
\end{abstract}

\vspace{0.5em}
\noindent\textbf{关键词：}货币政策执行报告；文本情绪；市场波动；结构断点；事件研究

\vspace{0.3em}
\noindent\textbf{JEL 分类号：}E52；E58；G14

% ---------------------------------------------------------------------
% 英文摘要（按需保留；若期刊或赛事不要求可删去）
% ---------------------------------------------------------------------
\clearpage
\renewcommand{\abstractname}{Abstract}
\begin{abstract}
\noindent
The Monetary Policy Report (MPR) published by the People's Bank of China is the
primary written channel through which policy intentions are communicated to the
market. Existing studies have established that the tone of the MPR explains
aggregate stock market returns, and that the forward-looking policy guidance
rather than the macroeconomic review drives this reaction. Most of this evidence,
however, rests on broad market indices, and whether it holds across institutional
periods has received little attention. Using 102 quarterly reports from 2001Q1 to
2026Q2 and a Chinese financial sentiment dictionary, we measure the tone of the two
segments separately. Market reactions are measured with event-window returns and
volatility constructed from an equal-weighted portfolio of 16 Shenwan level-1
industries.

We find no robust effect of policy guidance tone on the direction of post-release
returns, but a significant reduction in announcement-day volatility: a one
standard deviation increase in tone lowers the intraday range by 15 to 28 basis
points, depending on whether crisis years are included, and by 15 to 19 basis
points after controlling for pre-announcement volatility. The result survives an
alternative 28-industry market definition, a 500-draw permutation test, and
controls for reserve requirement ratio cuts, LPR cuts and reverse repo rate cuts,
while the tone of the macroeconomic review shows no such effect. The volatility
effect is concentrated before 2014; after 2014 the relation between tone and
returns turns negative, but this directional result is highly sensitive to crisis
years and is therefore reported as suggestive evidence.

These findings suggest that written central bank communication primarily reduces
disagreement rather than moving prices, and that its effect is not constant but
varies with the monetary policy framework. Treating heterogeneous periods and text
segments as a homogeneous sample may cancel out effects of opposite sign and lead
to the misleading conclusion that communication does not matter.
\end{abstract}

\vspace{0.5em}
\noindent\textbf{Keywords:} Monetary Policy Report; textual sentiment; market volatility;
structural break; event study

% ---------------------------------------------------------------------
% 写作依据（仅供团队核对，不进入排版）
%   "面向未来的政策指引而非宏观回顾"        → 姜富伟等（2021）的核心结论
%   前瞻性/回顾性分段的处理思路              → 王辉等（2023）；蒋海等（2025）
%   文本度量对词典与设定的敏感性（稳健性）    → 林建浩、孙乐轩（2025）
%   "信息效应"而非单纯冲击的分析框架          → 王熙、郑梦圆（2021）
%   行业事件窗与行业面板做法                  → 游庭玮（2026）；王春芝（2022）；周亮（2018）
% ---------------------------------------------------------------------
